\documentclass[11pt,a4paper]{article}
% =====================================================================
%  Shared preamble for the Part V lecture notes (lectures 19-22).
%
%  These four lectures sit outside Geron, so the notes ARE the primary
%  source and are examined on the same terms as the textbook parts.
%  Everything here is chosen to make that readable on paper: one column,
%  generous measure, and the same six-colour palette as the slides so a
%  student moving between the two is not re-learning what red means.
%
%  Build with pdflatex (twice, for the toc and the refs):
%      cd notes && latexmk -pdf lecture-19.tex
%  or  make -C notes
% =====================================================================
\usepackage[T1]{fontenc}
\usepackage[utf8]{inputenc}
\usepackage{newtxtext,newtxmath}      % Times-like: prints well, reads well
\usepackage[scaled=0.86]{inconsolata} % code, at a size that sits beside Times
\usepackage{microtype}
% newtx defines \Bbbk, \openbox, \proof and \endproof; amssymb and amsthm
% then redefine them and abort. Freeing the four names before those two
% packages load is the standard fix, and it costs nothing here: we use
% amsthm only for \newtheorem.
\let\Bbbk\relax
\usepackage{amsmath,amssymb}
\let\openbox\relax
\let\proof\relax
\let\endproof\relax
\usepackage{amsthm}
\usepackage{booktabs}
\usepackage{enumitem}
\usepackage{xcolor}
\usepackage{tcolorbox}
\usepackage{titlesec}
\usepackage{graphicx}
\usepackage[hidelinks]{hyperref}
\tcbuselibrary{skins,breakable}

% ---- the course palette, from assets/css/custom.css -------------------
\definecolor{aimlprimary}{HTML}{0B3D62}   % deep blue  - structure
\definecolor{aimlaccent} {HTML}{C0392B}   % brick red  - failures
\definecolor{aimlsuccess}{HTML}{14663A}   % green      - repairs
\definecolor{aimlmath}   {HTML}{6C3483}   % purple     - the derivation
\definecolor{aimlmuted}  {HTML}{4B5563}
\definecolor{aimlrule}   {HTML}{B0BCC7}
\definecolor{aimlsoft}   {HTML}{F4F7F9}

\usepackage[a4paper,margin=32mm,top=30mm,bottom=30mm]{geometry}
\setlength{\parskip}{0.45\baselineskip}
\setlength{\parindent}{0pt}
\linespread{1.06}

% ---- headings --------------------------------------------------------
\titleformat{\section}{\Large\bfseries\color{aimlprimary}}{\thesection}{0.7em}{}
\titleformat{\subsection}{\large\bfseries\color{aimlprimary}}{\thesubsection}{0.7em}{}
\titleformat{\subsubsection}{\bfseries\color{aimlmuted}}{}{0em}{}
\titlespacing*{\section}{0pt}{1.6\baselineskip}{0.5\baselineskip}
\titlespacing*{\subsection}{0pt}{1.2\baselineskip}{0.35\baselineskip}

% ---- boxes -----------------------------------------------------------
% One box per role, and the role is the colour. A student who has seen the
% slides already knows what each one means before reading a word of it.
\newtcolorbox{keybox}[1][]{
  breakable, enhanced, colback=aimlsoft, colframe=aimlprimary,
  boxrule=0.4pt, left=8pt, right=8pt, top=6pt, bottom=6pt,
  fonttitle=\bfseries\small, coltitle=white, colbacktitle=aimlprimary,
  attach boxed title to top left={yshift=-2mm, xshift=4mm},
  boxed title style={boxrule=0pt, arc=1pt}, #1}

\newtcolorbox{failbox}[1][]{
  breakable, enhanced, colback=white, colframe=aimlaccent,
  boxrule=0.9pt, left=8pt, right=8pt, top=6pt, bottom=6pt,
  fonttitle=\bfseries\small, coltitle=white, colbacktitle=aimlaccent,
  attach boxed title to top left={yshift=-2mm, xshift=4mm},
  boxed title style={boxrule=0pt, arc=1pt}, #1}

\newtcolorbox{derivbox}[1][]{
  breakable, enhanced, colback=white, colframe=aimlmath,
  boxrule=0.6pt, left=8pt, right=8pt, top=6pt, bottom=6pt,
  fonttitle=\bfseries\small, coltitle=white, colbacktitle=aimlmath,
  attach boxed title to top left={yshift=-2mm, xshift=4mm},
  boxed title style={boxrule=0pt, arc=1pt}, #1}

% An exercise. Deliberately the same shape as the notebook's `try` field:
% one change, and what should happen to the output.
\newtcolorbox{trybox}[1][]{
  breakable, enhanced, colback=white, colframe=aimlsuccess,
  boxrule=0.5pt, left=8pt, right=8pt, top=6pt, bottom=6pt,
  fonttitle=\bfseries\small, coltitle=white, colbacktitle=aimlsuccess,
  attach boxed title to top left={yshift=-2mm, xshift=4mm},
  boxed title style={boxrule=0pt, arc=1pt}, #1}

% ---- small semantics -------------------------------------------------
\newcommand{\scope}[1]{\textcolor{aimlmuted}{\small #1}}
\newcommand{\term}[1]{\textbf{#1}}
\newcommand{\code}[1]{\texttt{\small #1}}
\newcommand{\lecture}[1]{Lecture~#1}

\newtheorem{definition}{Definition}[section]
\newtheorem{proposition}[definition]{Proposition}

% ---- title block -----------------------------------------------------
% \notestitle{19}{Information retrieval: the lexical foundation}{SciFact (BEIR)}
\newcommand{\notestitle}[3]{%
  \thispagestyle{empty}
  \begin{flushleft}
    {\color{aimlmuted}\small\sffamily
     APPLICATIONS OF MACHINE LEARNING \quad$\cdot$\quad
     BSc Mathematics of Artificial Intelligence}\\[2mm]
    {\color{aimlrule}\rule{\textwidth}{0.8pt}}\\[4mm]
    {\color{aimlmuted}\large Lecture #1 \quad$\cdot$\quad Part V \quad$\cdot$\quad
     Extended notes}\\[3mm]
    {\Huge\bfseries\color{aimlprimary}\baselineskip=1.15\baselineskip #2\par}\vspace{4mm}
    {\color{aimlmuted}\large Dataset: #3}\\[3mm]
    {\color{aimlrule}\rule{\textwidth}{0.8pt}}
  \end{flushleft}
  \vspace{2mm}
  \begin{keybox}[title={Why these notes exist}]
    Lectures 19--22 sit outside G\'eron, the course's primary text. For those
    four lectures \emph{these notes are the primary source}, and they are
    examined on exactly the same terms as the textbook parts. Everything in
    the slide deck for this lecture is here, with the arguments written out at
    length rather than compressed onto a slide; the numbers are the ones the
    accompanying notebook prints, and every one of them is a measurement on
    the corpus named above rather than a figure quoted from a paper.
  \end{keybox}
  \vspace{1mm}
  {\small\tableofcontents}
  \vspace{2mm}
  \noindent{\color{aimlrule}\rule{\textwidth}{0.4pt}}
  \vspace{1mm}
}


\begin{document}
\notestitlefor{9}{II}{Neural networks, from the perceptron up}{Fashion-MNIST}{Chapter 9}

\section{The brief, and why this is not Lecture 3}

A document-processing bureau scans incoming items and routes each to one of ten
queues. Items arrive as small greyscale scans, already cropped and centred; a
human currently sorts them, and the operator wants to know how often a machine
would be right. Misrouting is recoverable but expensive --- the item goes to the
wrong desk and comes back a day later.

Ten classes, one label each, no ranking, no probability requirement stated.
Write that down: it decides the output layer.

\begin{center}
\small
\begin{tabular}{p{0.40\textwidth}p{0.44\textwidth}}
\toprule
\textbf{Lectures 3--4} & \textbf{Today} \\
\midrule
Two classes & ten classes \\
Wildly imbalanced & exactly balanced \\
Accuracy was a trap & accuracy is defensible --- we will check, not assume \\
The question was where to put the threshold & the question is what function to
                                              fit \\
\bottomrule
\end{tabular}
\end{center}

\begin{keybox}[title={The habit worth keeping}]
Not ``accuracy is bad''. \emph{Count the classes before you choose.} Here every
class holds exactly 6{,}000 images, so ``always $k$'' scores 10.0\% and
accuracy cannot hide a useless model.
\end{keybox}

\subsection{The corpus, and what it does not contain}

Fashion-MNIST: 60{,}000 training and 10{,}000 test images, $28\times28$
greyscale, ten balanced classes. We carve 5{,}000 validation images out of the
training half with a fixed seed --- 55{,}000 / 5{,}000 / 10{,}000 --- because
we are about to try ten configurations, and Lecture 2 established what happens
to a number chosen on the set it is reported on.

What the data does not contain: \emph{no colour}, so a red shirt and a blue
shirt are the same image; \emph{no scale}, since every garment fills the frame,
so size cannot distinguish a sandal from a boot; and \emph{no context}, one
centred item on a black background. Every one of those is a simplification the
bureau's real scans will not honour, so a number measured here is an upper
bound on what the same method gives on the archive --- and the report should
say so.

\begin{failbox}[title={Failure condition --- a pixel is a feature, and that costs something}]
We flatten each image into 784 numbers and hand it over as an ordinary table.
The model is told nothing about which pixels are neighbours: shuffle all 784
columns with one fixed permutation and it learns exactly as well, though a
human could not read a single shuffled image.

That is a real loss. Repairing it is four lectures away. Today we pay it, and
it is the honest starting point.
\end{failbox}

Dividing the pixels by 255 is not Lecture 2's leak. \code{StandardScaler}
estimates a mean and a variance from data and must be fitted after the split;
$x/255$ estimates nothing --- it divides by a constant that is a property of
the file format.

\section{Derivation: what a layer computes}

The vocabulary is historical. McCulloch and Pitts modelled a biological neuron
that fires when enough of its inputs fire, which is where \emph{neuron},
\emph{activation} and \emph{firing} come from. It is not an argument for
anything; modern networks resemble brains about as closely as aeroplanes
resemble birds.

A \term{threshold logic unit} computes
\[ z = \mathbf{w}^{\mathsf T}\mathbf{x} + b,
   \qquad \hat y = \operatorname{step}(z). \]
A weighted sum and a threshold --- and the weighted sum is Lecture 2's linear
model wearing a different name. Rosenblatt trained it with
\[ w_{i,j}^{\,\text{next}} = w_{i,j} + \eta\,(y_j - \hat y_j)\,x_i, \]
which converges only if the classes are linearly separable. \emph{No loss
function appears anywhere in that rule}, which is the tell: this is not
gradient descent, and it does not generalise.

And a single unit draws one straight line. XOR needs two, so no TLU computes
it --- Minsky and Papert said so in 1969 and the field stopped for fifteen
years. The repair is embarrassingly simple: stack them.

\subsection{One layer, in four steps}

\begin{derivbox}[title={From one unit to a mini-batch}]
A single unit computes $z = \mathbf{w}^{\mathsf T}\mathbf{x} + b$: two numbers
per unit and no more, a direction and how far along it to sit.

Stack $n$ units side by side and the weights become a matrix,
$\mathbf{z} = \mathbf{W}^{\mathsf T}\mathbf{x} + \mathbf{b}$ with one column
per unit. The layer is \emph{one matrix product}, which is why a tuned BLAS
matters at all.

Feed a whole mini-batch by making $\mathbf{X}$ one row per instance:
$\mathbf{Z} = \mathbf{X}\mathbf{W} + \mathbf{b}$, with shapes
$(m\times d)(d\times n) \to (m\times n)$ and $\mathbf{b}$ broadcasting along
the rows.

Apply a non-linearity elementwise:
\[ \mathbf{H} = \phi\bigl(\mathbf{X}\mathbf{W} + \mathbf{b}\bigr). \]
Elementwise, so it changes no shape --- the arithmetic above is untouched.
\end{derivbox}

\begin{derivbox}[title={Now remove $\phi$, and watch the stack collapse}]
Compose two layers with nothing between them:
\[ (\mathbf{X}\mathbf{W}_1 + \mathbf{b}_1)\mathbf{W}_2 + \mathbf{b}_2
   = \mathbf{X}(\mathbf{W}_1\mathbf{W}_2)
   + (\mathbf{b}_1\mathbf{W}_2 + \mathbf{b}_2). \]
Matrix product is associative, so the pair collapses to
$\mathbf{X}\mathbf{W}' + \mathbf{b}'$ --- one affine map. A stack of linear
layers, however deep, is a linear model.
\end{derivbox}

\begin{keybox}[title={So $\phi$ is not a refinement}]
Without it there is no second layer at all. Two hundred stacked layers would
still be a single affine map, and that is one line of algebra --- it is the
reason activation functions exist.
\end{keybox}

The step function had to go for a different reason: gradient descent needs a
slope, and the step's derivative is zero everywhere it exists.

\begin{center}
\small
\begin{tabular}{llll}
\toprule
\textbf{Activation} & \textbf{Range} & \textbf{Cost} & \textbf{Problem} \\
\midrule
Sigmoid & $(0,1)$ & an exponential & saturates at both ends \\
tanh    & $(-1,1)$ & an exponential & saturates at both ends \\
ReLU    & $[0,\infty)$ & a comparison & flat below zero \\
\bottomrule
\end{tabular}
\end{center}

$\max(0,z)$ is fast, does not saturate above, and in practice trains better.
``Flat below zero'' is a real defect with a name --- dying ReLUs --- diagnosed
in Lecture 11, not here.

\begin{keybox}[title={What the derivation buys you}]
A layer is one matrix product and one elementwise function; everything else is
bookkeeping. The parameter count is $d\times n + n$ per layer, so you can
\emph{price} a network before fitting it. Depth is only worth having because
$\phi$ is there. And the shapes compose, which is what makes a mini-batch free.

Lecture 10 differentiates exactly this expression, once per layer, backwards.
Lecture 11 asks what happens to the variance of $\mathbf{H}$ as the stack gets
deeper, and answers with the same two objects.
\end{keybox}

\section{From one unit to a network}

\subsection{Price it before you fit it}

\begin{center}
\small
\begin{verbatim}
784 x 300 + 300 = 235,500   layer 1
300 x 100 + 100 =  30,100   layer 2
100 x  10 +  10 =   1,010   output
                  -------
                  266,610   total
\end{verbatim}
\end{center}

266{,}610 parameters fitted from 55{,}000 images --- nearly five parameters per
training image. Lecture 2's failure condition for least squares was $m < n$,
more parameters than instances. We are past it comfortably, and the model will
still work. Understanding why is Lectures 13 and 14.

\subsection{The output layer, and the loss that must pair with it}

Ten mutually exclusive classes, so ten output units and a softmax across them:
\[ \hat p_k = \frac{\exp(z_k)}{\sum_{j=1}^{10}\exp(z_j)}. \]
Every output is positive and they sum to 1, and the largest logit gives the
largest probability, so \code{argmax} is unchanged by the softmax. The loss is
cross-entropy,
\[ L = -\sum_{k=1}^{10} y_k\log\hat p_k = -\log\hat p_{\text{true}}, \]
since the target is one-hot and nine of the ten terms are zero. Chapter 4
called this log loss when there were two classes; it is the same object.

\begin{center}
\small
\begin{tabular}{llll}
\toprule
\textbf{Task} & \textbf{Output units} & \textbf{Output activation} &
\textbf{Loss} \\
\midrule
Binary & 1 & sigmoid & binary cross-entropy \\
Multilabel binary & one per label & sigmoid & binary cross-entropy \\
Multiclass & one per class & softmax & cross-entropy \\
\bottomrule
\end{tabular}
\end{center}

Ours is the third row. The distinction between rows two and three is not
cosmetic: multilabel asks ten independent yes/no questions, multiclass asks one
ten-way question, and softmax couples the outputs where ten sigmoids do not.

Those tables answer three questions whose answers are \emph{forced} rather than
chosen --- how many output units (the shape of the answer), which activation
(the range the answer must live in), and which loss (what a wrong answer
costs).

\begin{failbox}[title={Failure condition --- the pairing is not optional}]
Softmax with mean squared error trains, slowly and badly: the gradient through
a saturated softmax nearly vanishes, and MSE does nothing to undo it.
Cross-entropy is the loss that cancels the softmax derivative, which Lecture 17
derives. Until then, take the pairing from the table.
\end{failbox}

Two names arrive here and are cashed later. \term{Backpropagation} is one
forward pass to compute the loss, one backward pass to compute every partial
derivative, and one descent step --- and the backward pass costs about as much
as the forward one \emph{no matter how many parameters there are}. 266{,}610
parameters, one number of interest, and a single sweep gets all the
derivatives; that claim is derived at the top of the next lecture and is the
whole reason any of this is feasible. And \term{Adam} keeps running averages of
the gradient and its square and scales each parameter's step by the second,
which makes it far less sensitive to the initial learning rate than plain SGD
--- though less sensitive is not insensitive, as the sweep below shows.

\section{Fitting it, and reading what happened}

\begin{center}
\small
\begin{tabular}{lr}
\toprule
\textbf{Quantity} & \textbf{Measured} \\
\midrule
Training images & 55{,}000 \\
Epochs & 20 \\
Wall clock, CPU & 125 s \\
Validation accuracy & 88.94\% \\
Test accuracy & 88.02\% \\
Baseline & 10.0\% \\
\bottomrule
\end{tabular}
\end{center}

Hold on to the wall clock --- it is the number the next lecture repairs.

Three words get confused, and Scikit-Learn does not help. A \term{batch} is the
instances used for one weight update (128 here); a \term{step} is one weight
update (430 per epoch); an \term{epoch} is one complete pass over the training
set (20 of them). Scikit-Learn calls the number of epochs \code{max\_iter}. It
is not iterations. Read the docstring, not the parameter name.

The learning curves separate and stay separated, and that gap is overfitting
--- Lecture 2's table drawn as lines. The loss falls smoothly, so the optimiser
is working; training accuracy keeps climbing, so the model has capacity to
spare; validation accuracy flattens, so the extra capacity is being spent on
memorising. And the curve answers a question the number cannot: would more
epochs help? Here, no --- the validation curve stopped moving before the budget
ran out, and without the curve you would have guessed.

\begin{center}
\small
\begin{tabular}{lll}
\toprule
\textbf{Repair for that gap} & \textbf{Available today} & \textbf{Taught in} \\
\midrule
Stop when validation stops improving & \code{early\_stopping=True} &
                                       Lecture 5, applied here \\
Penalise large weights & \code{alpha} & Lecture 5 \\
Dropout & not available & Lecture 10, then 11 \\
Batch normalisation & not available & Lecture 11 \\
More data, or augmented data & not available & Lecture 13 \\
\bottomrule
\end{tabular}
\end{center}

Three of the five are not reachable from this library at all --- and they all
need a hook inside the training loop.

\section{Ten fits, and three things they do not tell you}

Five architectures and five learning rates, on a 10{,}000-image subset so they
fit in the lecture.

\begin{center}
\small
\begin{tabular}{lrrr}
\toprule
\textbf{Hidden layers} & \textbf{Parameters} & \textbf{Validation accuracy} &
\textbf{Seconds} \\
\midrule
30            &  23{,}860 & 84.66\% & 169 \\
100           &  79{,}510 & 86.16\% & 225 \\
300           & 238{,}510 & 86.20\% & 347 \\
300, 100      & 266{,}610 & 84.80\% & 555 \\
300, 200, 100 & 316{,}810 & 84.24\% & 841 \\
\bottomrule
\end{tabular}
\end{center}

\begin{center}
\small
\begin{tabular}{lr}
\toprule
\textbf{Learning rate} & \textbf{Validation accuracy} \\
\midrule
0.0001 & 84.28\% \\
0.0003 & 86.22\% \\
0.001 --- the default & 84.80\% \\
0.003 & 85.76\% \\
0.01 & 84.20\% \\
\bottomrule
\end{tabular}
\end{center}

On this subset \emph{depth loses}: one hidden layer of 300 beats both deeper
networks, and the deepest is the worst and the slowest. And the library's
default learning rate is not the best value on the grid.

\begin{failbox}[title={Failure condition --- a sweep whose winner depends on the scale it was run at}]
The five architectures span about two accuracy points; so do the five learning
rates. The notebook sweeps 6{,}000 images rather than 10{,}000 and gets a
\emph{different} answer: there the learning rate spans 2.9 points and the
architectures 1.3, and even the winning rate changes.

Two runs of the same sweep, one lecture, two conclusions. Which is why ``the
learning rate matters more'' is a claim about a grid and a scale, never about
the world --- and why a sweep must always be reported with the size it was run
at.
\end{failbox}

Ten fits is also not a search: ten points from a space with five axes, varied
one axis at a time with everything else held fixed. You have not optimised
anything; you have looked at ten places.

\begin{keybox}[title={And the ranking does not transfer}]
\begin{center}
\small
\begin{tabular}{lrrl}
\toprule
\textbf{Run} & \textbf{Images} & \textbf{Epochs} & \textbf{Accuracy} \\
\midrule
Best in the sweep --- 300, at lr 0.001 & 10{,}000 & 12 & 86.20\% validation \\
The headline model --- 300, 100, at lr 0.001 & 55{,}000 & 20 & 88.02\% test \\
\bottomrule
\end{tabular}
\end{center}

The architecture the sweep ranked \emph{third of five} is the one that wins
with all the data. A deeper network needs more data to pay for its extra
parameters, so ranking architectures on a fifth of your data systematically
prefers the shallow ones --- and never finds out.
\end{keybox}

\section{One number over ten classes}

Split it. Lecture 3 built the confusion matrix for two classes and the object
generalises without change: row $i$, column $j$ counts the images of class $i$
the model called class $j$.

Recall on \emph{Shirt} is 74.1\%, against 98.1\% on \emph{Sneaker}. Shirts are
called Pullover 10.9\% of the time, T-shirt/top 7.8\%, and Coat 4.6\% ---
exactly the three classes anyone would name from the image grid, and Coat is
nearly as bad at 74.8\%. \emph{The model is not confused about anything a
person would find easy.}

\begin{keybox}[title={A defensible report}]
``Overall it is right 88.0\% of the time. But it is right 74.1\% of the time on
Shirt, and almost all of those errors go to three visually similar queues. If
those three desks are adjacent, the practical error rate is much lower than the
headline. If they are in different buildings, it is much worse.''

The headline accuracy does not answer the operator's question. The per-class
breakdown does, and it took one extra function call.
\end{keybox}

\section{Four walls, and they are all the same wall}

\begin{center}
\small
\begin{tabular}{p{0.30\textwidth}p{0.56\textwidth}}
\toprule
\textbf{What you want} & \textbf{What happens} \\
\midrule
Make it faster & there is no \code{device=}. Scikit-Learn does not support
                 GPUs and says so in its own FAQ --- a deliberate scope
                 decision, not an oversight \\
Change the loss & \code{loss="mae"} raises \code{TypeError}. Cross-entropy is
                 hard-coded \\
See a gradient & \code{[a for a in dir(clf) if "grad" in a.lower()]} is empty.
                 266{,}610 partial derivatives were computed on every batch and
                 none was kept \\
Stop part-way & \code{partial\_fit} is the finest handle there is, and one call
                is one complete pass over everything you hand it \\
\bottomrule
\end{tabular}
\end{center}

Each of these matters for an ordinary reason, not an exotic one. A weighted
loss is one line in a framework that lets you write one, and the operator
saying ``calling a coat a shirt is cheap, calling a bag a sandal is expensive''
is the normal case. Lecture 11 builds a twenty-layer network whose loss barely
moves, and the diagnosis is a per-layer gradient statistic you cannot compute
here.

\begin{failbox}[title={Failure condition --- the problem underneath all four}]
\code{fit()} contains a \code{for} loop over epochs and batches that you did
not write, cannot see, and cannot enter. Every one of the four walls is that
one wall.

This is not a criticism of Scikit-Learn --- \code{MLPClassifier} is documented
as a convenience, and for a tabular problem with a standard loss it is the
right tool. It is the wrong tool from here on, and now you know precisely why.
\end{failbox}

\section{The notebook}

\begin{trybox}[title={What to change}]
Remove the activation --- set it to identity --- and refit. Accuracy collapses
towards what a linear model gets, which is the collapse derived in \S2,
measured.
\end{trybox}

\section{Exercises}

\begin{enumerate}[leftmargin=1.6em,itemsep=4pt]
  \item Show that a stack of dense layers with identity activations is a single
        affine map, and state what that implies about depth. \hfill \scope{[5]}
  \item A network has 784 inputs, hidden layers of 300 and 100 units, and 10
        outputs. Compute the parameter count, and say which layer holds most of
        it and why. \hfill \scope{[4]}
  \item A task has ten mutually exclusive classes. Give the number of output
        units, the output activation and the loss, and explain why the same
        answer would be wrong for ten independent yes/no labels.
        \hfill \scope{[4]}
  \item A five-point architecture sweep on a fifth of the data ranks a deeper
        network third of five; trained on all the data it wins. Explain the
        mechanism, and say what it implies about tuning on subsamples.
        \hfill \scope{[5]}
  \item A model reports 88.0\% accuracy on ten balanced classes. State one
        further measurement you would make before reporting it to the
        stakeholder, and what decision it would inform. \hfill \scope{[4]}
\end{enumerate}

\section*{Summary}
\addcontentsline{toc}{section}{Summary}

A layer is one matrix product and one elementwise non-linearity,
$\mathbf{H} = \phi(\mathbf{X}\mathbf{W} + \mathbf{b})$, and without $\phi$ any
stack collapses to a single affine map by associativity --- so depth exists
only because the non-linearity does. The parameter count $d\times n + n$ per
layer lets you price a network before fitting it: 266{,}610 for ours, from
55{,}000 images.

The architecture tables answer three forced questions --- how many output
units, which activation, which loss --- and the loss must pair with the
activation rather than being chosen beside it. The network reached 88.02\% test
accuracy against a 10.0\% baseline in 125 seconds on a CPU.

Ten hand fits gave three lessons, none of them about the winner. Depth lost on
10{,}000 images; neither knob dominated \emph{on that grid at that scale}, and
the notebook's 6{,}000-image sweep reaches a different conclusion from the same
code; and the architecture ranked third of five is the one that wins with all
the data, because a deeper network needs more data to pay for its parameters.

Then four things \code{MLPClassifier} will not do --- run on a GPU, change the
loss, expose a gradient, or stop inside an epoch --- which are all the same
thing: the training loop was written for you, and you cannot enter it.

\end{document}
